% =====================================================================
%  Lettera di presentazione della candidatura
%  Va in: Candidatura/documenti/esterni/lettera-di-presentazione.tex
%  Sul sito compare in evidenza in cima alla Candidatura.
% =====================================================================
\documentclass[lettera]{gpt}

\titolo{Lettera di Presentazione}
\oggetto{Candidatura per il capitolato C2 -- NEXUM Concierge}
\data{2026-10-28}                  % data di invio della lettera
\luogo{Padova}
\responsabile{spagnoli}            % chi firma: il Responsabile
\destinatari{Prof.~Tullio Vardanega\newline
  Prof.~Riccardo Cardin\newline
  Università degli Studi di Padova\newline
  Dipartimento di Matematica ``Tullio Levi-Civita''\newline
  Via Trieste, 63 -- 35121 Padova}

% Link ai PDF nel repository: il nome non cambia con la versione, quindi
% il link porta sempre all'ultima versione approvata del documento.
\newcommand{\pdfdoc}[2]{\href{https://github.com/Git-Push-Team/docs/blob/main/Candidatura/#1}{#2}}

\begin{document}

Egregio Prof.~Vardanega, egregio Prof.~Cardin,

con la presente il gruppo \href{https://github.com/Git-Push-Team/}{\textbf{Git Push Team}} (gruppo~9) sottopone alla
Vostra attenzione la propria candidatura per la realizzazione del
capitolato \textbf{C2 -- NEXUM Concierge}, proposto da \textbf{Eggon S.r.l.}

La documentazione completa prodotta a supporto della candidatura è consultabile sul nostro sito
all'indirizzo \url{https://git-push-team.github.io/docs/}. Di seguito
riassumiamo le motivazioni della scelta, i costi e gli impegni assunti, e
riportiamo i collegamenti ai singoli documenti sul nostro repository "docs" dell'organizzazione \url{https://github.com/Git-Push-Team/docs}.

\subsection*{Motivazioni della scelta}

A seguito dell'analisi comparativa dei capitolati, il gruppo ha scelto C2
per la disponibilità dimostrata dall'azienda e l'equilibrio tra interesse, valore formativo e sostenibilità:
\begin{itemize}
  \item \textbf{perimetro circoscritto}: l'MVP è
        limitato al dominio di ferie e permessi, con vincoli espliciti su
        ciò che non è richiesto, e riduce il rischio di ampliamenti
        non pianificati dai requisiti;
  \item \textbf{criteri di accettazione verificabili}: le capacità
        richieste all'assistente conversazionale e il percorso completo, dalla domanda
        informativa alla registrazione della richiesta, sono una base
        solida su cui articolare requisiti e test;
  \item \textbf{tecnologie attuali in un prodotto reale}: l'assistente
        si integra nella piattaforma NEXUM e porta il
        gruppo a sperimentare approcci basati su LLM, Retrieval-Augmented Generation e MCP, elementi che hanno stimolato l'interesse di tutto il team;
  \item \textbf{proponente presente}: l'incontro con Eggon del 26 ottobre
        ha confermato la disponibilità di risorse cloud, un incontro
        settimanale o bisettimanale fisso e la volontà di un confronto proattivo sulle
        scelte progettuali.
\end{itemize}
In caso di mancata aggiudicazione, l'ordine di preferenza del gruppo è
C1 -- AURA (Vimar S.p.A.) e poi C0 -- AIRA (Sanmarco Informatica S.p.A.).
Il confronto completo è presente nel documento
\pdfdoc{documenti/esterni/analisi-dei-capitolati/analisi-dei-capitolati.pdf}{Analisi dei Capitolati}.

\subsection*{Costi e impegni}

Come indicato nella
\pdfdoc{documenti/esterni/dichiarazione-di-impegni/Dichiarazione\%20di\%20Impegni.pdf}{Dichiarazione di Impegni},
il gruppo preventiva un costo complessivo di \textbf{11.250,00~\texteuro},
per un totale di 552 ore di lavoro ripartite in modo uguale tra i sei
componenti (92 ore ciascuno), e si impegna a completare il progetto entro
il \textbf{23 marzo 2027}.

\subsection*{Documenti consultabili}

\begin{itemize}
  \item \pdfdoc{documenti/esterni/analisi-dei-capitolati/analisi-dei-capitolati.pdf}{Analisi dei Capitolati};
  \item \pdfdoc{documenti/esterni/dichiarazione-di-impegni/Dichiarazione\%20di\%20Impegni.pdf}{Dichiarazione di Impegni};
  \item verbali esterni:
    \begin{itemize}
      \item \pdfdoc{verbali/esterni/VE-2026-10-21.pdf}{21 ottobre 2026}, incontro con Sanmarco Informatica;
      \item \pdfdoc{verbali/esterni/VE-2026-10-22.pdf}{22 ottobre 2026}, incontro con Vimar;
      \item \pdfdoc{verbali/esterni/VE-2026-10-26.pdf}{26 ottobre 2026}, incontro con Eggon;
    \end{itemize}
  \item verbali interni:
    \begin{itemize}
      \item \pdfdoc{verbali/interni/VI-2026-10-12.pdf}{12 ottobre 2026}, costituzione del gruppo;
      \item \pdfdoc{verbali/interni/VI-2026-10-13.pdf}{13 ottobre 2026}, discussione e votazione dei capitolati;
      \item \pdfdoc{verbali/interni/VI-2026-10-19.pdf}{19 ottobre 2026}, contatti con i proponenti e assegnazione dei ruoli;
      \item \pdfdoc{verbali/interni/VI-2026-10-20.pdf}{20 ottobre 2026}, struttura del repository della documentazione;
      \item \pdfdoc{verbali/interni/VI-2026-10-26.pdf}{26 ottobre 2026}, discussione sui capitolati e azioni da intraprendere.
    \end{itemize}
\end{itemize}
Tutta la documentazione riportata è consultabile sul nostro \href{https://git-push-team.github.io/docs/}{sito} e \href{https://github.com/Git-Push-Team/docs}{repository}.

\subsection*{Componenti del gruppo}

\begin{center}
  \begin{tabular}{@{}lr@{}}
    \toprule
    \intestazione{Nominativo} & \intestazione{Matricola} \\
    \midrule
    \nome{biasini}    & 2147948 \\
    \nome{cescobolla} & 2138006 \\
    \nome{corro}      & 2145607 \\
    \nome{guglielmo}  & 2137975 \\
    \nome{mazzurana}  & 2149463 \\
    \nome{spagnoli}   & 2147977 \\
    \bottomrule
  \end{tabular}
\end{center}

Restiamo a disposizione per qualsiasi chiarimento.

\end{document}